\documentclass{article}
\usepackage[UTF8]{ctex}
\usepackage{amsmath,amssymb,amsfonts}
\usepackage{enumitem}
\usepackage{geometry}
\geometry{a4paper, margin=1in}
% === 错题本排版优化（不参与题目标识） ===
\setlength{\parindent}{0pt}
\setlength{\parskip}{0.35em plus 0.10em minus 0.08em}
\setlength{\emergencystretch}{3em}
\setlist{itemsep=0.25em, topsep=0.40em, parsep=0pt}
\allowdisplaybreaks[2]
% === 错题本排版优化结束 ===
\begin{document}
\section*{北京科技大学2025年高等代数真题}
\section*{10．北京科技大学2025年高等代数试卷}

\subsection{题目1}
1．（15 分）设齐次线性方程组

\[
\left\{\begin{array}{c}
(1+a) x_{1}+2 x_{2}+3 x_{3}+\cdots+n x_{n}=0 \\
x_{1}+(2+a) x_{2}+3 x_{3}+\cdots+n x_{n}=0 \\
\cdots \cdots \\
x_{1}+2 x_{2}+3 x_{3}+\cdots+(n+a) x_{n}=0
\end{array}\right.
\]

其中 \(n \geq 2\) 。问：\(a\) 取何值时，该方程组有非零解，并求其通解。

\subsection{题目2}
2．（20 分）已知 \(A\) 为 \(n\) 阶可逆的反对称矩阵，\(b\) 为 \(n\) 元列向量，若矩阵 \(B=\left(\begin{array}{cc}A & b \\ b^{\mathrm{T}} & 0\end{array}\right)\) ，求矩阵 \(B\) 的秩．

\subsection{题目3}
3．（15 分）设 \(A=\left(a_{i j}\right)_{n \times n}\) ，其元素 \(a_{i j}\) 都为整数 \((i, j=1,2, \cdots, n)\) 。令

\[
d=\left|\begin{array}{ccccc}
a_{11}-\frac{1}{k} & a_{12} & a_{13} & \cdots & a_{1 n} \\
a_{21} & a_{22}-\frac{1}{k} & a_{23} & \cdots & a_{2 n} \\
\vdots & \vdots & \vdots & & \vdots \\
a_{n 1} & a_{n 2} & a_{n 3} & \cdots & a_{n n}-\frac{1}{k}
\end{array}\right| .
\]

这里 \(k\) 为正整数且 \(k \geq 2\) 。证明：\(d \neq 0\) 。

\subsection{题目4}
4．（15 分）设 \(M_{3}(\mathbb{R})\) 为所有3阶实方阵按矩阵的加法及实数与矩阵的数量乘法构成的实数域 \(\mathbb{R}\) 上的线性空间．已知 \(M_{3}(\mathbb{R})\) 的两个子空间

\[
W_{1}=\left\{\left.\left(\begin{array}{ccc}
a & 0 & c \\
a & 0 & 0 \\
c & b & 0
\end{array}\right) \right\rvert\, a, b, c \in \mathbb{R}\right\}, W_{2}=\left\{\left.\left(\begin{array}{ccc}
x & 0 & 0 \\
0 & y & 0 \\
0 & 0 & z
\end{array}\right) \right\rvert\, x, y, z \in \mathbb{R}\right\} .
\]

\begin{itemize}
  \item[（1）] 求和空间 \(W_{1}+W_{2}\) 的维数和一组基。
  \item[（2）] 记 \(W=W_{1}+W_{2}\) ，求子空间 \(W_{3}\) ，使得 \(M_{3}(\mathbb{R})=W_{3} \oplus W\) ，并说明理由。
\end{itemize}

\subsection{题目5}
5．（15 分）设 \(V\) 是数域 \(P\) 上的 \(n\) 维线性空间，\(\alpha_{1}, \alpha_{2}, \cdots, \alpha_{n}\) 为 \(V\) 的一组基，设 \(\sigma\) 为 \(V\) 上的线性变换，且 \(\sigma\) 在 \(\alpha_{1}, \alpha_{2}, \cdots, \alpha_{n}\) 下的矩阵为 \(A\) ，记齐次线性方程组 \(A X=0\) 的基础解系为 \(\eta_{1}, \eta_{2}, \cdots, \eta_{n-r}\) ，其中\(r\) 为矩阵 \(A\) 的秩，令 \(\beta_{i}=\left(\alpha_{1}, \alpha_{2}, \cdots, \alpha_{n}\right) \eta_{i}(i=1,2, \cdots, n-r)\) ，证明： \(\operatorname{Ker} \sigma=L\left(\beta_{1}, \beta_{2}, \cdots, \beta_{n-r}\right)\) 。

\subsection{题目6}
6．（15 分）设 \(B\) 为正定矩阵，\(A-B\) 为半正定矩阵。证明：

\begin{itemize}
  \item[（1）] 满足 \(|A-2 \lambda B|=0\) 的所有 \(\lambda\) 都不小于 \(\frac{1}{2}\) ．
  \item[（2）] \(|A| \geq|B|\) ．
\end{itemize}

\subsection{题目7}
7．（15 分）设多项式 \(f(x)=\left(x-a_{1}\right)\left(x-a_{2}\right)\left(x-a_{3}\right)\left(x-a_{4}\right)-1\) ，其中 \(a_{1}, a_{2}, a_{3}, a_{4}\) 是 4 个互不相同的整数．问：\(f(x)\) 在有理数域上是否可约？

\subsection{题目8}
8．（20 分）设 \(V\) 是数域 \(P\) 上的 \(n\) 维线性空间，\(n \geq 2\) ，设 \(\sigma\) 为 \(V\) 上的线性变换，若有 \(\xi \in V\) ，满足\(\sigma^{n-1}(\xi) \neq 0, \sigma^{n}(\xi)=0\) ，求 \(V\) 的一组基，使得 \(\sigma\) 在这组基下的矩阵为

\[
\left(\begin{array}{ccccc}
0 & 0 & \cdots & 0 & 0 \\
1 & 0 & \cdots & 0 & 0 \\
0 & 1 & \cdots & 0 & 0 \\
\vdots & \vdots & & \vdots & \vdots \\
0 & 0 & \cdots & 1 & 0
\end{array}\right)
\]

\subsection{题目9}
9．（20 分）设 \(A\) 为 \(n\) 阶正定矩阵，其中 \(n \geq 2\) ，若 \(\alpha\) 为任意的 \(n\) 维非零列向量，\(B=A \alpha \alpha^{\mathrm{T}}\) 。

\begin{itemize}
  \item[（1）] 证明：\(\left|\lambda E_{n}-B\right|=\lambda^{n-1}\left(\lambda-\alpha^{\mathrm{T}} A \alpha\right), E_{n}\) 为 \(n\) 阶单位矩阵。
  \item[（2）] 求 \(B\) 最大的特征值 \(\lambda_{\text {max }}\) ．
  \item[（3）] 求 \(B\) 的属于 \(\lambda_{\text {max }}\) 的特征子空间的维数与一组基。
\end{itemize}
\end{document}
